\documentclass[12pt,oneside]{book}

% ============================================================
% METADATOS
% ============================================================
\newcommand{\universidad}{Universidad de Buenos Aires (UBA)}
\newcommand{\facultad}{Facultad de Derecho}
\newcommand{\materia}{Derecho Procesal Civil y Comercial}
\newcommand{\titulo}{Organización y Competencia Judicial}
\newcommand{\autor}{Isaac Edgar Camacho Ocampo}
\newcommand{\anio}{2025}

% ============================================================
% PAQUETES
% ============================================================
\usepackage[utf8]{inputenc}
\usepackage[T1]{fontenc}
\usepackage[spanish,es-nodecimaldot]{babel}

\usepackage[
    top=2.5cm,
    bottom=2.5cm,
    left=3cm,
    right=2.5cm,
    headheight=15pt
]{geometry}

\usepackage{amsmath}
\usepackage{amssymb}
\usepackage{mathtools}

\usepackage{graphicx}
\usepackage{float}
\usepackage{caption}
\usepackage{subcaption}

\usepackage{booktabs}
\usepackage{array}
\usepackage{longtable}
\usepackage{tabularx}

\usepackage{enumitem}

\usepackage{xcolor}
\usepackage[most]{tcolorbox}

\usepackage{fancyhdr}
\usepackage{ragged2e}

% ============================================================
% RUTAS DE IMÁGENES
% ============================================================
\graphicspath{
    {./img/}
    {./graficos/}
    {./figuras/}
}

% ============================================================
% CONFIGURACIÓN GENERAL
% ============================================================
\setcounter{tocdepth}{2}

\setlength{\parindent}{0pt}
\setlength{\parskip}{0.5em}

\setlist[itemize]{
    topsep=0.3em,
    itemsep=0.2em
}

\setlist[enumerate]{
    topsep=0.3em,
    itemsep=0.2em
}

\clubpenalty=10000
\widowpenalty=10000

% ============================================================
% ENCABEZADOS Y PIES
% ============================================================
\pagestyle{fancy}
\fancyhf{}

\fancyhead[L]{\nouppercase{\leftmark}}
\fancyhead[R]{\materia}

\fancyfoot[C]{\thepage}

\renewcommand{\headrulewidth}{0.4pt}
\renewcommand{\footrulewidth}{0pt}

\fancypagestyle{plain}{
    \fancyhf{}
    \fancyfoot[C]{\thepage}
    \renewcommand{\headrulewidth}{0pt}
}

% ============================================================
% COMANDOS MATEMÁTICOS
% ============================================================
\newcommand{\abs}[1]{\left|#1\right|}
\newcommand{\norm}[1]{\left\lVert#1\right\rVert}
\newcommand{\vect}[1]{\mathbf{#1}}
\newcommand{\deriv}[2]{\frac{d#1}{d#2}}
\newcommand{\pderiv}[2]{\frac{\partial#1}{\partial#2}}

% ============================================================
% CAJAS ACADÉMICAS
% ============================================================
\newtcolorbox{definition}[1]{
    enhanced,
    breakable,
    title={Definición: #1},
    fonttitle=\bfseries,
    colback=blue!3,
    colframe=blue!50!black,
    boxrule=0.8pt,
    arc=2mm
}

\newtcolorbox{important}{
    enhanced,
    breakable,
    title={Importante},
    fonttitle=\bfseries,
    colback=yellow!8,
    colframe=orange!70!black,
    boxrule=0.8pt,
    arc=2mm
}

\newtcolorbox{example}[1]{
    enhanced,
    breakable,
    title={Ejemplo: #1},
    fonttitle=\bfseries,
    colback=green!4,
    colframe=green!40!black,
    boxrule=0.8pt,
    arc=2mm
}

\newtcolorbox{formula}[1]{
    enhanced,
    breakable,
    title={Fórmula / Regla: #1},
    fonttitle=\bfseries,
    colback=purple!4,
    colframe=purple!50!black,
    boxrule=0.8pt,
    arc=2mm
}

\newtcolorbox{exercise}[1]{
    enhanced,
    breakable,
    title={Ejercicio: #1},
    fonttitle=\bfseries,
    colback=gray!5,
    colframe=gray!60!black,
    boxrule=0.8pt,
    arc=2mm
}

\newtcolorbox{note}{
    enhanced,
    breakable,
    title={Nota},
    fonttitle=\bfseries,
    colback=white,
    colframe=gray!50,
    boxrule=0.6pt,
    arc=2mm
}

% ============================================================
% HIPERVÍNCULOS Y METADATOS PDF
% ============================================================
\usepackage[
    colorlinks=true,
    linkcolor=blue!50!black,
    citecolor=green!40!black,
    urlcolor=blue!60!black,
    pdftitle={\titulo},
    pdfauthor={\autor},
    pdfsubject={Apuntes universitarios},
    bookmarks=true
]{hyperref}

% ============================================================
% DOCUMENTO
% ============================================================
\begin{document}

% ------------------------------------------------------------
% PORTADA
% ------------------------------------------------------------
\begin{titlepage}

\thispagestyle{empty}

\begin{center}

\vspace*{1cm}

{\large \textsc{\universidad}}

\vspace{0.2cm}

{\large \textsc{\facultad}}

\vspace{3cm}

{\Huge \textbf{\materia}}

\vspace{1cm}

{\LARGE \titulo}

\vspace{0.8cm}

\textit{Comprender el porqué de cada norma es el primer paso hacia su correcta aplicación.}

\vspace{1.2cm}

\rule{0.70\textwidth}{0.5pt}

\vspace{0.5cm}

{\large Apuntes de estudio}

\vspace{0.5cm}

\rule{0.70\textwidth}{0.5pt}

\vfill

{\large \autor}

\vspace{0.3cm}

{\large \anio}

\end{center}

\vfill

\begin{flushleft}
\textit{Este es un modesto aporte a los estudiantes de ``Derecho Procesal Civil y Comercial'' no pretende sustituir el material oficial de la materia.}
\end{flushleft}

\end{titlepage}

% ------------------------------------------------------------
% ÍNDICE
% ------------------------------------------------------------
\tableofcontents

% ------------------------------------------------------------
% CONTENIDO
% ------------------------------------------------------------

\chapter{Organización y Competencia Judicial}

En esta unidad se desarrollan los conceptos fundamentales de la competencia judicial dentro de la teoría general del derecho procesal. Se parte de la ubicación de la competencia en relación con los conceptos de proceso, jurisdicción y acción, explicando cómo el Estado distribuye los conflictos entre los distintos órganos jurisdiccionales. Se analizan los caracteres esenciales de la competencia (orden público, improrrogabilidad e indelegabilidad), sus excepciones en materia patrimonial y territorial, y los distintos criterios de distribución: fuero ordinario y federal, materia, territorio, valor, grado y turno. Se aborda además la organización del sistema judicial argentino (doble estructura federal y ordinaria), la competencia de la Corte Suprema de Justicia de la Nación (originaria y apelada), y los mecanismos procesales para discutir la competencia: la declinatoria y la inhibitoria, con sus posibles derivaciones en contiendas negativas o positivas.

\begin{note}
Una analogía útil para relacionar los tres conceptos centrales del curso es pensar en una ciudad con muchos médicos especialistas. La jurisdicción equivale al sistema de salud pública en su conjunto: el Estado que resuelve los conflictos. La competencia equivale a la especialidad de cada médico, ya que no cualquier profesional puede operar del corazón. El proceso, por su parte, equivale a los protocolos y pasos que sigue cada intervención. Así, la competencia organiza qué juez atiende qué conflicto, del mismo modo en que la especialidad médica determina quién trata cada enfermedad.
\end{note}

\section{Ubicación de la Competencia en la Teoría General del Derecho Procesal}

La competencia debe ubicarse dentro de la teoría general del derecho procesal y vincularse con los demás temas centrales de la materia: el proceso, la jurisdicción y la acción.

En el lenguaje coloquial suele utilizarse equívocamente la palabra ``competencia'' en lugar de ``jurisdicción''. Sin embargo, todos los jueces están investidos con el mismo poder jurisdiccional: no existe mayor jurisdicción en un juez de la Corte que en un juez de primera instancia, sino una distribución distinta de la tarea.

\begin{quote}
\itshape
``No quiero decir con esto que hay mayor jurisdicción en un juez de la Corte que en un juez de primera instancia, sino que lo que hay es una distribución distinta de la tarea.''
\end{quote}

A mayor concentración demográfica, mayor especialización de los jueces.

\begin{table}[H]
\centering
\caption{Conceptos centrales de la teoría general del proceso}
\begin{tabularx}{\textwidth}{>{\raggedright\arraybackslash}p{0.18\textwidth} X >{\raggedright\arraybackslash}p{0.22\textwidth}}
\toprule
\textbf{Concepto} & \textbf{Definición} & \textbf{Rol en el sistema} \\
\midrule
Jurisdicción & Mecanismo por el cual el Estado quita a los ciudadanos la posibilidad de resolver conflictos usando la fuerza (teoría de la sustitución). & El poder estatal \\
Acción & La posibilidad de acceder al sistema jurisdiccional para resolver el conflicto. & El acceso al sistema \\
Proceso & La regla de juego que fija los límites de intervención de la jurisdicción. & La regla de juego \\
Competencia & La distribución material de los distintos conflictos entre los distintos jueces: la medida de la jurisdicción. & Quién resuelve qué conflicto \\
\bottomrule
\end{tabularx}
\end{table}

\begin{definition}{Competencia}
La competencia es la distribución material de los distintos conflictos entre los distintos jueces. Constituye la medida de la jurisdicción: mientras la jurisdicción es un poder único e indivisible del Estado, la competencia determina qué porción de ese poder ejerce cada juez en concreto.
\end{definition}

\begin{example}{Especialización territorial según densidad poblacional}
\begin{itemize}
\item \textbf{Ciudad de Buenos Aires:} jueces sumamente especializados (jueces de familia, jueces civiles patrimoniales, jueces comerciales).
\item \textbf{Provincia de Buenos Aires (interior):} juzgados civiles y comerciales que se ocupan de todos esos temas conjuntamente.
\item \textbf{Sur del país (Patagonia):} tribunales con competencia universal, en los que el juez conoce en absolutamente todos los conflictos del ámbito territorial, dado que hay poca población y menos conflictos.
\end{itemize}
\end{example}

\section{Caracteres de la Competencia}

\subsection{Orden Público: la regla general}

La primera regla relativa a la competencia es que se trata de una cuestión de orden público. Esto significa que, cuando el conflicto se suscita, las personas no pueden llevarlo a cualquier juez: ya existe una norma del Estado que, previamente al conflicto, determina quién es el juez competente. Esto responde a la garantía constitucional del juez natural.

\begin{important}
Existe una garantía constitucional —la del juez natural del conflicto— conforme a la cual la norma que determina quién es el juez competente es previa a que el conflicto se suscite.
\end{important}

\begin{quote}
\itshape
``Hay una garantía constitucional que es la del juez natural del conflicto: esa norma estaría ya previamente a que el conflicto se suscite determinando quién es el juez que se tiene que ocupar de entender en ese conflicto.''
\end{quote}

\subsection{Improrrogabilidad: con sus excepciones}

\begin{formula}{Artículo 1º del CPCCN}
``La competencia atribuida a los tribunales nacionales es improrrogable.'' Esto significa que no puede ser pactada por las partes que están en el conflicto.
\end{formula}

Este principio tiene una importante excepción que suele ser fuente de confusión.

\begin{important}
Solo se puede prorrogar la competencia cuando se reúnen dos condiciones en forma simultánea:
\begin{enumerate}
\item Que se trate de materia \textbf{patrimonial} (no de familia, no penal, no laboral, no federal).
\item Que la prórroga sea respecto de la competencia \textbf{territorial}.
\end{enumerate}
La prórroga puede ser \textbf{expresa} (pactada en un contrato) o \textbf{tácita} (el demandado contesta sin cuestionar la competencia territorial).
\end{important}

\begin{quote}
\itshape
``Si estamos en una cuestión patrimonial, sí se puede prorrogar. No se podría hacer, por ejemplo, en una cuestión de familia, o en una cuestión penal, o en una cuestión laboral. La materia no se puede pactar, así como tampoco se puede pactar la esfera federal.''
\end{quote}

\begin{example}{Contrato de alquiler con cláusula de prórroga}
En un contrato de alquiler, las partes pueden establecer: ``si surgiere algún conflicto relativo a este contrato, lo llevaremos a un juzgado de la Ciudad de Buenos Aires, o de la Provincia de Buenos Aires, o de la Provincia de Santa Fe''. Esto es válido porque: (a) es materia patrimonial, y (b) se está pactando sobre competencia territorial.
\end{example}

Sobre la prórroga tácita:

\begin{quote}
\itshape
``El actor lo va a llevar al juez que quiera que sea el competente en función del territorio, sin haberse establecido ningún tipo de pacto previo. Y la otra parte —la demandada— si al contestar la demanda no dice nada, está prorrogando tácitamente esa atribución, cuando podría por las reglas generales llevarlo a algún otro magistrado.''
\end{quote}

Las reglas generales a las que se refiere son, por ejemplo, el lugar donde sucedió el hecho o el domicilio del demandado (regla fijada en el CPCCN).

\subsection{Indelegabilidad}

El tercer carácter esencial es la indelegabilidad: el juez al que se le asigna una competencia determinada no puede excusarse de entender en el caso. No puede derivarlo a un colega de igual grado ni a la segunda instancia.

\begin{quote}
\itshape
``Es un deber para el juez intervenir en todos los asuntos en los que esté asignada a su competencia.''
\end{quote}

Esto no significa que el juez acepte sin más toda atribución de competencia que las partes le hagan: de hecho, puede y debe controlar su competencia. En la jurisprudencia de la Corte Nacional, la mayor cantidad de fallos se refieren precisamente a discusiones de competencia, lo que refleja que las normas no siempre son claras y los mismos conflictos pueden interpretarse de uno u otro modo.

\begin{note}
Los litigantes también eligen ciertos fueros para llevar sus conflictos porque saben que son fueros más proclives a acceder al reclamo que van a hacer, lo cual genera adicionalmente discusión sobre quién es el juez competente para cada caso.
\end{note}

\section{Organización del Sistema Judicial Argentino: Doble Estructura}

En Argentina conviven dos estructuras paralelas de organización judicial. Cuando en 1853-1860 se conforma el Estado argentino, lo hace a partir de la existencia de estados previos —las provincias— que delegaron ciertas facultades en el Estado federal y retuvieron las restantes. La materia procesal fue entendida como una función no delegada, con lo cual cada provincia se dio su propio código procesal y organizó su propia justicia local.

\begin{quote}
\itshape
``La regulación de los procesos —todo lo que es materia procesal— a diferencia de lo que es el derecho sustancial, se ha entendido —no sin algún tipo de discusión en su momento— que era una función que los estados preexistentes no delegaban en el estado federal. Hoy se concluyó que es una función no delegada, con lo cual cada estado se va a dar su código procesal y se va a dar su organización en cuanto a lo que es la justicia local u ordinaria.''
\end{quote}

\begin{table}[H]
\centering
\caption{Doble estructura del sistema judicial argentino}
\begin{tabularx}{\textwidth}{X X}
\toprule
\textbf{Jurisdicción Federal} & \textbf{Jurisdicción Ordinaria / Local} \\
\midrule
Presente en todo el territorio nacional & Organizada por cada provincia (poder no delegado) \\
Órgano máximo: Corte Suprema de Justicia de la Nación & Órgano máximo: superiores tribunales de cada provincia \\
Código aplicable: CPCCN y leyes federales & Código aplicable: código procesal propio de cada provincia \\
\bottomrule
\end{tabularx}
\end{table}

\subsection{Los Tribunales Nacionales: una situación histórica particular}

Los tribunales nacionales con asiento en la Ciudad de Buenos Aires dependen funcionalmente de la esfera federal (Corte Suprema, Consejo de la Magistratura de la Nación), pero no conocen en materia federal, sino en asuntos locales ordinarios (accidentes de tránsito, relaciones laborales, cuestiones penales no federales, etc.).

\begin{note}
\textbf{Contexto histórico: el territorio cedido.} La Ciudad de Buenos Aires (antes denominada Capital Federal) era territorio de la Provincia de Buenos Aires. Cuando esta ingresó a la Federación en el Pacto de San José de Flores, cedió ese territorio al gobierno federal para que se asentaran las autoridades nacionales (Congreso, Poder Ejecutivo, Poder Judicial).

Dato de color: el Banco de la Provincia de Buenos Aires —institución preexistente a la Federación— es territorio provincial dentro de la Capital Federal. Quien entra al banco está pisando la Provincia de Buenos Aires dentro de la Capital Federal.
\end{note}

En 1853-1860, además de la Capital, había numerosos territorios nacionales (la Patagonia, Tierra del Fuego, el norte) aún no provincializados. La Ciudad de Buenos Aires, densamente poblada, necesitaba tribunales propios. Como era territorio cedido por la provincia, fue el gobierno federal quien organizó esos tribunales para que atendieran los asuntos de los ciudadanos de la ciudad. Esos son los actuales tribunales nacionales, que dependen de la Corte federal y del Consejo de la Magistratura de la Nación.

\begin{quote}
\itshape
``Un juez nacional para ser designado tiene que concursar en el Consejo de la Magistratura de la Nación y es designado con acuerdo del Senado del Congreso Nacional y por un decreto del Presidente de la República. Sin embargo, estos jueces no conocen en asuntos o en materia federal, sino que se ocupan de asuntos locales.''
\end{quote}

En 1994, con la reforma constitucional, se creó la Ciudad Autónoma de Buenos Aires (CABA), que adquirió autonomía propia (prácticamente la de una provincia). En 1996, la ciudad sancionó su propio estatuto (su constitución). Desde entonces, lo lógico hubiera sido que los tribunales nacionales pasaran a depender de la esfera local. Sin embargo, esto no ocurrió de manera completa:

\begin{quote}
\itshape
``Desde el 94-96 hasta la actualidad hay una disputa y solo se han pasado ciertas competencias penales a la esfera de la Ciudad de Buenos Aires. Pero el núcleo duro de lo que son los juzgados nacionales —a nivel penal, laboral, civil y comercial— permanece en la esfera federal.''
\end{quote}

\section{Criterios de Distribución de la Competencia}

El legislador utiliza distintos criterios para distribuir los conflictos entre los diferentes jueces. Estos criterios se superponen: en cada conflicto concreto puede interactuar más de uno para determinar el juez competente.

\subsection{Criterio 1: Fuero Federal vs. Ordinario (Local)}

Es la primera gran diferenciación. Ciertos conflictos le interesan al Estado federal y se resolverán en el fuero federal, mientras que otros quedan en la órbita local y se resolverán ante la justicia ordinaria de cada provincia o de la Ciudad de Buenos Aires.

\subsection{Criterio 2: La Materia}

Más allá de la distinción federal-ordinario, el tipo de conflicto determina la especialidad del tribunal.

\begin{table}[H]
\centering
\caption{Fueros ordinarios típicos}
\begin{tabularx}{\textwidth}{>{\raggedright\arraybackslash}p{0.32\textwidth} X}
\toprule
\textbf{Fuero} & \textbf{Tipo de conflicto} \\
\midrule
Civil y Comercial & Conflictos patrimoniales: contratos \\
Laboral & Relaciones de trabajo y empleados \\
Familia & Divorcio \\
Penal (no federal) & Delitos del Código Penal sin carácter federal \\
Contencioso-administrativo & Conflictos con el Estado local \\
\bottomrule
\end{tabularx}
\end{table}

\subsection{Criterio 3: El Territorio}

El territorio también distribuye los conflictos: se prefiere que intervenga el juez más próximo al lugar del conflicto y las partes.

\begin{example}{Competencia territorial en el fuero de familia}
En cuestiones de familia se asigna como competencia territorial el último lugar donde convivió esa familia: la sede del hogar. Ese lugar estaría fijando la proximidad entre el tribunal y el conflicto. Así, ese tribunal con competencia territorial es el que entiende en los conflictos que se susciten con posterioridad a la separación: el régimen de comunicación con los hijos, los alimentos, etc.
\end{example}

\subsection{Criterio 4: El Valor o Monto}

El valor en disputa también puede determinar competencias. Las causas de menor cuantía pueden asignarse a tribunales específicos (en muchas provincias existen juzgados de pequeñas causas). En el ámbito del CPCCN, el criterio de valor opera a través del monto de apelabilidad.

\begin{formula}{Artículo 242 del CPCCN}
Las causas que tengan un monto inferior al que determina la Corte Nacional no van a tener la posibilidad de ser apeladas.
\end{formula}

\begin{note}
El monto vigente fue modificado a partir del comienzo de 2025. Al momento de la clase se indicó que era de \$300.000 (pesos argentinos).
\end{note}

\subsection{Criterio 5: El Grado}

% TODO: contenido incompleto o ambiguo en las notas originales — el desarrollo específico de este criterio no consta con mayor detalle en el apunte original más allá de su enunciación como criterio de distribución vinculado con las instancias de revisión de los conflictos.

\subsection{Criterio 6: El Turno}

El turno distribuye los asuntos que se van suscitando en un momento determinado del tiempo. Se aplica principalmente en materia penal y de familia, donde los conflictos pueden surgir en cualquier momento.

\begin{quote}
\itshape
``Hay tribunales que están abiertos las 24 horas para recibir todas las denuncias que se hacen en ese espacio temporal, y eso quiere decir que en ese turno todas esas causas van a ser asignadas a ese tribunal, que se va a ocupar de prevenir en la causa y después de quedarse en el conocimiento definitivo de esa discusión.''
\end{quote}

\section{Competencia de la Corte Suprema de Justicia de la Nación}

\subsection{Competencia Originaria}

\begin{formula}{Artículo 117 de la Constitución Nacional}
Regula la competencia originaria de la Corte Suprema. Esta interviene como tribunal de primera y única instancia en los conflictos:
\begin{itemize}
\item Que involucren diplomáticos, cónsules extranjeros o algún Estado extranjero en conflicto con un Estado local o ciudadano local.
\item Conflictos entre dos provincias argentinas, o entre un ciudadano de una provincia y otra provincia.
\end{itemize}
\end{formula}

El fundamento de esta atribución es evitar conflictos internos entre provincias, o situaciones que puedan comprometer las relaciones exteriores de la Argentina.

\subsection{Política Restrictiva de la Corte: el caso Barreto}

Desde hace más de diez años, la Corte ha desarrollado una política de restricción de su competencia, tanto en la instancia originaria como en la apelada.

\begin{table}[H]
\centering
\caption{Modelos comparados de tribunales supremos}
\begin{tabularx}{\textwidth}{>{\raggedright\arraybackslash}p{0.38\textwidth} X}
\toprule
\textbf{Modelo} & \textbf{Descripción} \\
\midrule
Modelo 1 — Alta cantidad de causas & La Corte interviene en muchísimas causas, como históricamente fue el caso argentino. \\
Modelo 2 — Menor cantidad, mayor relevancia & El tribunal superior interviene en menos casos pero con mayor trascendencia institucional. Es el modelo de la Corte Suprema de los Estados Unidos. \\
\bottomrule
\end{tabularx}
\end{table}

\begin{quote}
\itshape
``La Corte argentina decía —y seguimos teniendo— una Corte con un sinnúmero de incumbencias. Ya desde el 2007 en adelante empezó a restringir de alguna manera la cantidad de casos en los que iba a intervenir.''
\end{quote}

\begin{example}{Caso Barreto y el concepto de ``causa civil''}
A partir del caso Barreto se definió el concepto de ``causa civil'': solo se interviene en instancia originaria cuando la cuestión de derecho sustancial es central en la decisión del caso, no cuando es meramente tangencial. Cuando la discusión requiere análisis de normas locales o administrativas provinciales, la Corte indica que la causa debe iniciarse ante un tribunal federal de primera instancia con asiento en ese estado.
\end{example}

\begin{quote}
\itshape
``Es una línea bastante difusa para trazar, pero en los casos en que la discusión requiere, por ejemplo, el análisis de normas locales o de normas administrativas de ese estado provincial, la Corte viene diciendo desde entonces que, por aplicación del precedente Barreto, corresponde que la causa no se inicie en instancia originaria de la Corte, sino que deberá interponerse ante un tribunal federal de primera instancia con asiento en ese estado.''
\end{quote}

\subsection{Competencia Apelada}

La Corte tiene dos tipos de competencia apelada.

\paragraph{a) Competencia ordinaria apelada.} Era procedente en asuntos donde fuera parte el Estado federal y se superara un determinado monto.

\begin{example}{Caso Anadón}
La Corte Suprema declaró inconstitucional la competencia ordinaria apelada del Estado, clausurando la posibilidad de que el Estado nacional llegue a la Corte como tercera instancia ordinaria de revisión en causas donde fuera parte, con el fundamento de que el Estado no debía tener el privilegio de llegar a una tercera (o cuarta) instancia de revisión.
\end{example}

\paragraph{b) Competencia extraordinaria apelada.} Se conforma principalmente por el Recurso Extraordinario Federal (REF), regulado en el artículo 14 de la Ley 48, cuando se suscita una cuestión federal. A esto se suman otros supuestos de creación jurisprudencial de la propia Corte.

\begin{table}[H]
\centering
\caption{Vías de acceso a la competencia apelada extraordinaria}
\begin{tabularx}{\textwidth}{>{\raggedright\arraybackslash}p{0.28\textwidth} X}
\toprule
\textbf{Vía} & \textbf{Descripción} \\
\midrule
REF (Art. 14, Ley 48) & Recurso Extraordinario Federal. Procede cuando se suscita una cuestión federal. Vía ordinaria de acceso a la Corte en su competencia apelada. \\
Arbitrariedad & Supuesto de creación jurisprudencial: procede cuando la sentencia del tribunal inferior es arbitraria (carece de fundamento o es irrazonablemente fundada). \\
Gravedad Institucional & Casos de singular relevancia para las instituciones del país. También de creación jurisprudencial. \\
Per Saltum & Originariamente creación jurisprudencial; luego recibió sanción legislativa. Permite saltear instancias intermedias y pasar de primera instancia directamente a la Corte cuando la índole de la cuestión requiere intervención rápida y no admite demoras. Son casos muy excepcionales. \\
\bottomrule
\end{tabularx}
\end{table}

\begin{formula}{Artículo 280 del CPCCN — el ``certiorari'' negativo argentino}
En línea con la política restrictiva de competencias, la Corte aplica el artículo 280 del CPCCN para desestimar recursos extraordinarios sin más, cuando considera que los casos no alcanzan el umbral de trascendencia requerido. Este mecanismo, denominado ``certiorari negativo'', le permite seleccionar los casos que realmente va a revisar.
\end{formula}

\section{Cómo se Discute la Competencia: Mecanismos Procesales}

\subsection{Control de Oficio por el Juez}

\begin{quote}
\itshape
``La primera función que tiene el juez cuando recibe un caso es analizar si es competente. Y si no es competente, debe declarar su incompetencia y remitirla al tribunal que considere competente.''
\end{quote}

Existe una excepción a esta regla: si se trata de materia patrimonial y la incompetencia es territorial —que es prorrogable— el juez no puede declararse incompetente de entrada. Debe esperar a que el demandado interponga una excepción o cuestione la competencia. Si el demandado no lo hace, habrá prórroga tácita.

En todos los demás casos (competencia improrrogable), el juez controla de oficio y, si se considera competente, continúa con el trámite (corre traslado de la demanda); si se considera incompetente, lo declara mediante resolución fundada. Esa resolución puede ser apelada por las partes.

\subsection{La Declinatoria}

\begin{formula}{Artículo 8º del CPCCN}
Regula la declinatoria como una de las vías que tiene el demandado para cuestionar la competencia del juez que ya tiene la causa.
\end{formula}

La declinatoria es una excepción procesal. El demandado se presenta ante el mismo juez que está entendiendo en la causa y le plantea la incompetencia.

\begin{quote}
\itshape
``El demandado puede interponer una diversidad de excepciones que están previstas a partir del artículo 346-347 [del CPCCN], entre ellas la excepción de incompetencia: el demandado, cuando contesta la demanda ante el juez que tiene la causa, le dice: `Señor juez, usted es incompetente, pido que así lo declare.' Esto se discute con la otra parte y el juez deberá hacer una resolución donde acepte la excepción y resuelva enviar la causa al juez competente, o la rechace.''
\end{quote}

\subsection{La Inhibitoria}

La inhibitoria es la vía alternativa: en lugar de plantear la incompetencia ante el juez que tiene la causa, el demandado acude al juez que considera competente y le pide que se declare competente y solicite al otro juez que se inhiba (que deje de conocer en la causa).

\begin{quote}
\itshape
``El demandado que se notifica de la demanda en vez de ir al juez que tiene la causa, va al juez que él considera competente, le pide que se declare competente y que pida al juez que conoce en esa causa que se inhiba —de ahí el nombre de inhibitoria— de seguir conociendo esa causa.''
\end{quote}

La inhibitoria es mucho menos frecuente en la práctica:

\begin{quote}
\itshape
``Es menos común. De hecho, he visto pocos casos en ya 25 años de ejercicio.''
\end{quote}

\begin{example}{Conflictos entre tribunales nacionales y de CABA}
A raíz del conflicto de competencias entre tribunales nacionales y tribunales de la Ciudad de Buenos Aires, los jueces nacionales rechazaban la asignación de competencia a los tribunales locales. Los abogados del gobierno de la Ciudad optaron por utilizar directamente la vía inhibitoria: iban a un juez de la Ciudad con la causa ya iniciada ante un juez nacional y le pedían que se declarara competente y que solicitara la inhibición del juez nacional.
\end{example}

\begin{important}
La declinatoria y la inhibitoria son vías alternativas y excluyentes: si se usa una, no se puede usar la otra. Y en ambos casos, el demandado debe también contestar la demanda en el mismo acto, porque en el CPCCN las excepciones se oponen junto con la contestación de la demanda.
\end{important}

\begin{quote}
\itshape
``Aunque se use la inhibitoria, la presentación debe ser completa, aunque se haga ante el juez que no tiene la causa.''
\end{quote}

\subsection{Contiendas de Competencia: Positivas y Negativas}

Una vez planteada la cuestión (por declinatoria o inhibitoria), los jueces pueden estar o no de acuerdo.

\begin{table}[H]
\centering
\caption{Situaciones posibles tras plantear la cuestión de competencia}
\begin{tabularx}{\textwidth}{>{\raggedright\arraybackslash}p{0.28\textwidth} X}
\toprule
\textbf{Situación} & \textbf{Consecuencia} \\
\midrule
Acuerdo entre los jueces & Queda asignada la competencia. Una de las partes puede apelar, pero el caso sigue. \\
Contienda positiva & Ambos jueces se consideran competentes y quieren retener la causa. Debe intervenir el tribunal superior común para resolver. \\
Contienda negativa & Ambos jueces se consideran incompetentes y ninguno quiere asumir la causa. Debe intervenir el tribunal superior común. \\
\bottomrule
\end{tabularx}
\end{table}

\subsection{Tribunal que Resuelve las Contiendas}

La regla es que resuelve el tribunal superior común de los tribunales entre los cuales se da la contienda.

\begin{note}
Para evitar que todas las contiendas entre tribunales nacionales (todos pertenecientes a la órbita federal) lleguen a la Corte, una acordada establece que el tribunal que resuelve esas contiendas —sean positivas o negativas— es la Cámara del tribunal que previno en el caso.
\end{note}

Existe un caso donde inevitablemente interviene la Corte: cuando la contienda se da entre un tribunal de la Ciudad Autónoma de Buenos Aires y un tribunal nacional. Pese a que ambos tienen asiento en la ciudad, uno pertenece a la órbita de la justicia local de CABA y el otro a la justicia federal, con lo cual no tienen tribunal superior común salvo la Corte.

\begin{quote}
\itshape
``Inevitablemente la cuestión va a terminar en la Corte nacional. Y esto es el motivo por el cual si uno revisa la jurisprudencia de la Corte va a encontrar muchísimos fallos referidos a cuestiones de competencia: tanto más fallos relativos a estas cuestiones que fallos de cuestiones sustanciales o de derecho sustancial.''
\end{quote}

\section{Reflexión Final}

\begin{quote}
\itshape
``Lo importante en este curso de grado es saber y conocer las normas fundamentales y las reglas que sí o sí van a estar presentes, como son los caracteres tanto de la competencia en sí misma como los caracteres de asignación, y luego la organización de los tribunales. Porque después, cuando el caso concreto se suscite, uno inevitablemente —son casos que se dan todo el tiempo, pero cuando hay ciertos casos que no son tan típicos— inevitablemente el operador tiene que ponerse a estudiar para ver en qué tribunal corresponde que recaiga esa causa. Y además porque hay cuestiones que son opinables, que no están definidas de una vez y para siempre, sino que hay vaivenes en la jurisprudencia sobre cuál es el tribunal que corresponde entender en cada caso.''
\end{quote}

% ------------------------------------------------------------
% CORNELL
% ------------------------------------------------------------
\section{Cuadro Cornell — Síntesis de la Unidad}

\begin{table}[H]
\centering
\begin{tabularx}{\textwidth}{
    >{\raggedright\arraybackslash}p{0.28\textwidth}
    >{\raggedright\arraybackslash}X
}
\toprule
\textbf{Preguntas / palabras clave} &
\textbf{Notas / respuestas} \\
\midrule

\textbf{¿Qué es la competencia y cómo se diferencia de la jurisdicción?} &
La jurisdicción es el poder del Estado de resolver conflictos (poder único e indivisible). La competencia es la distribución de ese poder entre los distintos jueces: la medida de la jurisdicción. No hay mayor jurisdicción en un juez de la Corte que en uno de primera instancia; lo que varía es la distribución de la tarea. \\

\textbf{¿Cuáles son los tres caracteres esenciales de la competencia?} &
Orden público (norma previa que establece quién es el juez competente antes del conflicto, en garantía del juez natural); improrrogabilidad (como regla, las partes no pueden pactar el juez, salvo excepción); indelegabilidad (el juez no puede excusarse de entender en los asuntos asignados a su competencia). \\

\textbf{¿Cuándo es válida la prórroga de competencia?} &
Solo cuando se cumplen dos condiciones simultáneas: (a) que sea materia patrimonial, y (b) que la prórroga sea respecto de la competencia territorial. Puede ser expresa o tácita. No es posible prorrogar fueros de familia, penal, laboral ni la esfera federal. \\

\textbf{¿Cómo se organiza el sistema judicial argentino?} &
Existe una doble estructura: jurisdicción federal (presente en todo el territorio, con la Corte Suprema como cabeza) y jurisdicción ordinaria o local (organizada por cada provincia, con sus superiores tribunales), como consecuencia de que la materia procesal es una facultad no delegada al Estado federal. \\

\textbf{¿Cuál es la situación particular de los tribunales nacionales de CABA?} &
Fueron organizados históricamente por el gobierno federal cuando la Capital era territorio cedido por la Provincia de Buenos Aires. Dependen funcionalmente de la Corte federal, pero conocen en asuntos ordinarios locales. Desde 1994, con la autonomía de CABA, el traspaso a la órbita local solo se ha completado parcialmente en materia penal. \\

\textbf{¿Cuáles son los criterios de distribución de la competencia?} &
Los criterios se superponen en cada conflicto: fuero (federal vs. ordinario), materia, territorio, valor, grado y turno. \\

\textbf{¿Qué es la competencia originaria de la Corte? ¿Qué cambió con Barreto?} &
La competencia originaria (Art. 117 CN) hace que la Corte actúe como tribunal de primera y única instancia en casos que involucren diplomáticos, cónsules o Estados extranjeros, y en conflictos entre provincias. Desde el caso Barreto, la Corte restringió su competencia originaria en ``causas civiles'': solo interviene cuando el derecho sustancial federal es central. \\

\textbf{¿Qué ocurrió con la competencia ordinaria apelada de la Corte?} &
Fue declarada inconstitucional en el caso Anadón, por considerar que el Estado no debe tener el privilegio de acceder a una tercera instancia de revisión. La competencia apelada de la Corte se reduce básicamente al REF y a los casos de arbitrariedad, gravedad institucional y per saltum. \\

\textbf{¿Qué es la declinatoria y cómo funciona?} &
Es una excepción procesal (Art. 8 CPCCN) mediante la cual el demandado cuestiona la competencia del juez que tiene la causa presentándose ante ese mismo juez, junto con la contestación de la demanda. \\

\textbf{¿Qué es la inhibitoria y en qué se diferencia de la declinatoria?} &
Son vías alternativas y excluyentes. En la inhibitoria, el demandado acude al juez que él considera competente y le pide que solicite al juez que tiene la causa que se inhiba de seguir conociendo. \\

\textbf{¿Qué son las contiendas de competencia y quién las resuelve?} &
Si los jueces no coinciden en la atribución de competencia, se genera una contienda: positiva (ambos quieren ser competentes) o negativa (ninguno quiere serlo). La resuelve el tribunal superior común; entre tribunales nacionales, la Cámara del tribunal que previno; si la contienda es entre un tribunal de CABA y uno nacional, interviene inevitablemente la Corte. \\

\midrule

\multicolumn{2}{p{\textwidth}}{
\textbf{Resumen:}
La competencia es el mecanismo mediante el cual el Estado distribuye los conflictos entre los distintos órganos jurisdiccionales. Se ubica en la teoría general junto con la jurisdicción y la acción. Sus caracteres esenciales son el orden público, la improrrogabilidad (con la excepción de materia patrimonial y territorial) y la indelegabilidad. El sistema judicial argentino presenta una doble estructura federal-ordinaria, derivada del federalismo constitucional. Los criterios de distribución de competencia se superponen: fuero, materia, territorio, valor, grado y turno. La Corte Suprema tiene competencia originaria (Art. 117 CN) y apelada extraordinaria (REF, arbitrariedad, gravedad institucional, per saltum), habiendo restringido progresivamente su intervención desde 2007 (casos Barreto y Anadón). Para discutir la competencia, el demandado cuenta con dos vías alternativas y excluyentes: la declinatoria y la inhibitoria. Los desacuerdos entre jueces generan contiendas positivas o negativas, que resuelve el tribunal superior común; las contiendas entre tribunales nacionales y de CABA desembocan inevitablemente en la Corte.
} \\

\bottomrule
\end{tabularx}
\end{table}

\end{document}
